\documentclass[11pt,a4paper]{article}
\usepackage[margin=2.0cm]{geometry}
\usepackage{amsmath,amssymb,booktabs,enumitem,hyperref,xcolor}
\usepackage[T1]{fontenc}
\setlist{itemsep=2pt,topsep=4pt}
\newcommand{\E}{\mathbb{E}}
\title{SwarmPlan: a multi-stage network-design model\\ for bottom-up electrification pathways}
\author{Working formulation -- prototype in \texttt{swarmplan/SwarmPlan.jl}}
\date{16 November 2026}
\begin{document}
\maketitle

\noindent\textbf{Purpose.} Decide, for one village, \emph{what to build, where, and when} as demand
unfolds: nanogrid-level solar and storage, peer-to-peer (swarm) lines between nanogrids, and
an optional upgrade to a mini-grid hub. The model extends the equilibrium model of Ouanes,
Gonz\'alez Grand\'on, Fagerli and Gabriel (2026, \emph{Energy Economics} 163, 109598) in four
directions: space, price-responsive demand, staged investment under uncertainty, and the
swarm-to-mini-grid upgrade.

\section{Cookbook: the model in six steps}
\begin{enumerate}
  \item \textbf{Scenario tree.} Nodes $n\in\mathcal{N}$ with parent $a(n)$, probability $\pi_n$,
        demand growth factor $\gamma_n$. Stages have equal length.
  \item \textbf{Sites.} Nanogrids $i\in\mathcal{I}$ with coordinates; index $0$ is the hub at the
        village centroid. Candidate swarm lines $e=(i,j)\in\mathcal{E}$ (nearest neighbours) with
        length $\ell_e$; feeders hub--site with length $\ell_i^{0}$.
  \item \textbf{Time.} Representative days $r$ with weights $\omega_r$ (sum $=7$ days), hours $h$.
  \item \textbf{Demand.} Load $D_{irh}$ is split into blocks $k$ with share $\sigma_k$ and
        willingness to pay $w_k$ (decreasing). Served demand $d_{ikrhn}\le\sigma_k\gamma_nD_{irh}$.
  \item \textbf{Investment.} Capacities and lines are cumulative and irreversible along the tree.
  \item \textbf{Objective.} Maximise expected welfare $=$ consumer value $-$ cost.
\end{enumerate}

\section{Decision variables}
\begin{tabular}{@{}ll@{}}
\toprule
$S_{in},B_{in}\ge0$ & PV (kWp) and battery (kWh) at site $i$, node $n$; $i=0$ is the hub\\
$y_{en}\in\{0,1\}$ & swarm line $e$ built\\
$z_{in}\in\{0,1\}$ & feeder hub--site $i$ built\\
$u_n\in\{0,1\}$ & hub built (mini-grid upgrade)\\
$g,c^{+},c^{-},s\ge0$ & solar used, charge, discharge, state of charge (per $i,r,h,n$)\\
$d_{ikrhn}\ge0$ & served demand in block $k$\\
$f_{e}^{\rightarrow},f_{e}^{\leftarrow}\ge0$ & swarm flows; $\phi_i^{\downarrow},\phi_i^{\uparrow}\ge0$ feeder flows\\
\bottomrule
\end{tabular}

\section{Formulation}
\begin{align}
\max\quad & \sum_{n}\frac{\pi_n}{T}\Big[\underbrace{\sum_{i,k,r,h}\omega_r w_k d_{ikrhn}}_{\text{consumer value}}
 - \underbrace{\sum_{i\ge1}\big(c^{S}S_{in}+c^{B}B_{in}\big)}_{\text{nanogrid assets}}
 - \underbrace{\theta\big(c^{S}S_{0n}+c^{B}B_{0n}\big)-F u_n}_{\text{hub}} \notag\\
 &\qquad - \sum_{e}c^{L}\ell_e y_{en}-\sum_i c^{F}\ell^{0}_i z_{in} - \text{OPEX}_n\Big]
 \qquad(\theta<1\text{: scale economies},\ F\text{: fixed cost}) \label{obj}
\end{align}
subject to, for all $n$ (and $i,r,h$ where they appear):
\begin{align}
& g_{irhn}+c^{-}_{irhn}-c^{+}_{irhn}
  +\!\!\sum_{e\ni i}\!\big[(1-\lambda\ell_e)f^{\text{in}}_{e}-f^{\text{out}}_{e}\big]
  +(1-\lambda^{F}\ell^{0}_i)\phi^{\downarrow}_{i}-\phi^{\uparrow}_{i}=\sum_k d_{ikrhn}
  && \text{site balance}\label{bal}\\
& g_{0rhn}+c^{-}_{0rhn}-c^{+}_{0rhn}+\sum_i\big[(1-\lambda^{F}\ell^{0}_i)\phi^{\uparrow}_{i}-\phi^{\downarrow}_{i}\big]=0
  && \text{hub balance}\\
& s_{irhn}=(1-\delta)s_{ir,h-1,n}+\eta c^{+}_{irhn}-c^{-}_{irhn}/\eta,\quad s_{ir0n}=s_{irHn}
  && \text{storage}\\
& g_{irhn}\le S_{in}\,\rho_{rh},\quad c^{\pm}_{irhn}\le\kappa B_{in},\quad s_{irhn}\le B_{in}
  && \text{limits}\\
& f_e\le\bar f\,y_{en},\qquad \phi_i\le\bar\phi\,z_{in},\qquad z_{in}\le u_n,\qquad S_{0n},B_{0n}\le M u_n
  && \text{network}\\
& d_{ikrhn}\le\sigma_k\gamma_nD_{irh}
  && \text{demand}\\
& c^{S}S_{in}+c^{B}B_{in}\le m_i\,\gamma_n\quad(i\ge1)
  && \text{budget}\label{bud}\\
& S_{in}\ge S_{i,a(n)},\ B_{in}\ge B_{i,a(n)},\ y_{en}\ge y_{e,a(n)},\ z_{in}\ge z_{i,a(n)},\ u_n\ge u_{a(n)}
  && \text{irreversible}\label{irr}
\end{align}
Non-anticipativity is built in by indexing decisions on tree nodes.

\paragraph{Compatibility (optional).} A share $\chi\in[0,1]$ of nanogrid-level hardware becomes
unusable once the hub exists. Usable capacities $\tilde S_{in}\le S_{in}$ and
$\tilde B_{in}\le B_{in}$ replace $S,B$ in the capacity limits for $i\ge1$, with
\[
\tilde S_{in}\le(1-\chi)S_{in}+\bar S(1-u_n),\qquad \tilde B_{in}\le(1-\chi)B_{in}+\bar B(1-u_n).
\]
Costs are still paid on $S,B$, so stranded hardware is a sunk cost.

\section{Pathways as restrictions of one model}
\begin{tabular}{@{}lll@{}}
\toprule
Pathway & Restriction & Interpretation\\
\midrule
Isolated & $y=z=u=0$ & solar home systems / stand-alone nanogrids\\
Swarm & $u=0$ & interconnect, never centralise\\
Mini-grid upfront & $u_n=1\ \forall n$, $S_{in}=B_{in}=0$ for $i\ge1$ & classic mini-grid\\
Staged upgrade & $u$ fixed to a pattern, e.g.\ $(0,0,1)$ & centralise only if demand is high\\
Ladder & none & best of all\\
\bottomrule
\end{tabular}

\medskip\noindent Because every pathway is a restriction of the ladder, differences in the optimal value
are directly interpretable: $V^{\text{ladder}}-\max\{V^{\text{swarm}},V^{\text{mini-grid}}\}$ is the
\textbf{option value of modularity}, the quantity the swarm literature asserts but has not measured.

\section{Operator and regulation layer}
A single operator owns all assets and posts one tariff $p$. Customers buy block $k$ iff
$w_k\ge p$. The operator solves \eqref{bal}--\eqref{irr} without \eqref{bud} and with objective
$\max\; p\sum d - \text{cost}$. With step demand the optimal tariff lies in $\{w_k\}$, so the
operator problem is solved by enumeration; a price cap $\bar p$ restricts the set to $w_k\le\bar p$.
Welfare is evaluated ex post with the consumer value in \eqref{obj}.

\paragraph{Next step (the OR contribution).} Replace enumeration by a bilevel program:
\[
\max_{\bar p,\ \text{subsidy}}\ \text{welfare}\big(x^\ast\big)\quad\text{s.t.}\quad
x^\ast\in\arg\max_x\{\text{operator profit}(x;\bar p,\text{subsidy})\},\ \ \text{profit}\ge0 ,
\]
with continuous inverse demand $p_i(q)=\alpha_i-\beta_iq$. The lower level is then a convex QP whose
KKT system gives an MPEC (Gabriel et al.\ 2012, ch.\ 6); with several competing operators it becomes
an EPEC. Binary network decisions sit in the upper level or are handled by decomposition.

\section{How this differs from Ouanes et al.\ (2026)}
\begin{tabular}{@{}p{3.2cm}p{5.6cm}p{6.6cm}@{}}
\toprule
 & Paper & SwarmPlan\\
\midrule
Demand & fixed & price-responsive (WTP blocks)\\
Welfare & weighted score, converted at an assumed rate & consumer value minus cost, in EUR\\
Interconnection & free, single bus & lines chosen, cost per metre, losses\\
Time & one week, one design & staged on a scenario tree\\
Mini-grid & not modelled & hub upgrade with fixed cost and scale economies\\
Market power & changes transfers only & changes quantities and investment\\
Size & 3 nanogrids & 8 sites (prototype), any number in principle\\
\bottomrule
\end{tabular}

\section{What is assumed, not measured}
Line costs, capacities and losses; hub fixed cost $F$ and scale factor $\theta$; WTP levels and
shares; site positions and budgets beyond the paper's three values; demand growth factors. Load and
solar profiles are synthetic stand-ins calibrated to the paper's figures. Results from the prototype
show mechanisms and orders of magnitude, not estimates.

\section{Known limits and extensions}
\begin{itemize}
  \item \textbf{Budget constraint.} \eqref{bud} ties asset location to ownership. With one operator the
        real friction is household credit, so the budget should be on upfront payments, with an
        operator cost of capital. Interconnection then competes with pay-as-you-go finance.
  \item \textbf{Solver.} Fixed-charge network design has a weak LP relaxation. HiGHS needs minutes for
        the two-stage case; Gurobi, or Benders decomposition by tree node, is the route to larger
        villages and deeper trees.
  \item \textbf{Lead times.} Investments are available in the node where they are decided.
  \item \textbf{Grid arrival.} Add a node attribute ``grid arrives'' with a feeder to the main grid;
        this links to the undergrid mini-grid literature.
  \item \textbf{National scale.} Fit a reduced form (welfare gain as a function of settlement size,
        spacing and budget inequality) and use it as a fourth technology in OnSSET.
\end{itemize}
\end{document}
